%% Shared content for the academic CV and industry resume.
%% Edit section text here; both documents load these definitions.
%% Reorder sections in cv.tex or resume.tex, then rebuild with latexmk.
%% Use \cvonly{...} or \resumeonly{...} for variant-specific content.
%% Compile cv.tex or resume.tex; this file contains section definitions only.

%% ------------------------------------------------------------ header
\newcommand{\HeaderSection}{%
  \header{Yile (Wayne) Wang}{%
    Dallas, TX\sep +1 (682) 248-6331\sep
    \href{mailto:ylwwayne@gmail.com}{\ul{ylwwayne@gmail.com}}\sep
    \link{https://github.com/yilewang}{GitHub}\sep
    \link{https://yilewang.github.io/}{Website}\sep
    \link{https://www.linkedin.com/in/yile-wang-02a070131/}{LinkedIn}\sep
    \link{https://scholar.google.com/citations?user=bgnVj8wAAAAJ}{Google Scholar}}}

%% ------------------------------------------------------------ summary
\newcommand{\SummarySection}{%
  \cvsection{Summary}
  \resumeonly{%
  Computational neuroscientist who designs and ships AI agent systems.
  Creator and maintainer of \link{https://github.com/yilewang/llm-for-zotero}{LLM-for-Zotero} (\textbf{3,100+ GitHub stars, 650K+ plugin downloads, 100+ releases}), a TypeScript research agent with an adaptive agent runtime, multi-provider model layer, context budgeting, and verified, undoable actions.
  Research on brain-aligned deep networks (\venue{Communications Biology}) and on how representational drift shapes neural decoding (first-author ICLR 2027 submission).}
  \cvonly{%
  Ph.D. candidate in computational neuroscience at UT Dallas.
  I study how neural representations drift over time and what that means for decoding, and how human neural data can shape artificial vision models.
  Co-developed ReAlnet, an EEG-aligned vision model (\venue{Communications Biology}, 2026); first-author work on representational drift under review at ICLR 2027 and accepted at CCN 2026.
  Creator of \link{https://github.com/yilewang/llm-for-zotero}{LLM-for-Zotero}, an open-source AI research agent with \textbf{3,100+ GitHub stars} and \textbf{650K+ plugin downloads}.}}

%% ------------------------------------------------------------ education
\newcommand{\EducationSection}{%
  \cvsection{Education}
  \leftright{\textbf{The University of Texas at Dallas}}{\textbf{Richardson, TX}}
  \leftright{\indented{Ph.D., Cognition and Neuroscience (GPA: 3.94/4)}}{Aug 2021 -- Aug 2027 (expected)}
  \leftright{\indented{M.S., Applied Cognition and Neuroscience, Computational Modeling track (GPA: 3.93/4)}}{Aug 2019 -- May 2021}
  \entrygap
  \leftright{\textbf{Southern Medical University}}{\textbf{Guangzhou, China}}
  \leftright{\indented{B.S., Psychology}}{Aug 2014 -- May 2018}}

%% ------------------------------------------------------------ experience
\newcommand{\ExperienceSection}{%
  \cvsection{Research \& Work Experience}
  \leftright{\textbf{Doctoral Researcher}, The University of Texas at Dallas}{Aug 2021 -- Present}
  \subentry{Representational drift: biologically constrained neural network modeling (with Kaushik Lakshminarasimhan)}
  \begin{itemize}
    \item Designed and trained neural networks to explain representational drift, translating biological constraints into model architecture, learning rules, and training strategies.
    \item Adapted learning rules, model parameters, and training procedures to capture biological mechanisms while meeting modeling objectives.
    \item Developed an efficient-coding account of representational drift; first-author paper under review at \textbf{ICLR 2027}.
    \item Built a biologically plausible \textbf{downstream decoding model} that uses \textbf{online statistical learning} to compensate for representational drift. Validated the model on the longitudinal THINGS-EEG2 dataset; manuscript \textbf{under review}.
  \end{itemize}
  \subentry{Brain-aligned vision models, ReAlnet (with Zitong Lu and Julie Golomb, Ohio State University)}
  \begin{itemize}
    \item Co-designed a multi-layer image-to-brain encoding framework that jointly trains ImageNet-pretrained CORnet-S for object classification and EEG generation. Combined mean-squared-error (MSE) and contrastive losses for EEG generation, yielding \textbf{10 individualized models} aligned to each subject's EEG.
    \item ReAlnets outperformed the \textbf{baseline model, CORnet-S}, in EEG representational similarity across all four visual layers (\textbf{up to 40\% relative improvement at peak}). Gains extended to unseen fMRI data and human behavioral benchmarks on Brain-Score.
    \item Built the \textbf{PyTorch training and validation pipeline}, including Brain-Score behavioral evaluation and representational similarity analysis. Ran control experiments (\textbf{loss ablations, unpaired and scrambled EEG}) isolating the contribution of neural alignment.
    \item Extended the framework to \textbf{ResNet-18}, achieving significant improvements in EEG representational similarity and demonstrating applicability beyond the original architecture.
  \end{itemize}
  \indented{\textbf{Publications:} \venue{Communications Biology} (2026; earlier version at an ICLR 2024 workshop) and \venue{Cognitive Neurodynamics} (2025).}\par
  \subentry{Computational modeling and biomarker discovery in Alzheimer's disease}
  \begin{itemize}
    \item Predicted conversion from mild cognitive impairment (MCI) to Alzheimer's disease \textbf{4 years before clinical diagnosis with 80\% specificity}. Used computational biomarkers from individualized brain-network models and multivariate discriminant analysis.
    \item Built \textbf{individualized digital brain twins in The Virtual Brain}, using dynamical systems constrained by structural MRI and diffusion imaging. Applied mean-field modeling and criticality analysis to characterize collective brain dynamics in MCI and Alzheimer's disease.
    \item Used \textbf{Python/pandas} to clean and integrate \textbf{high dimensional, multimodal data from hundreds of patients}, spanning MRI, cognitive assessments, and behavioral measures. Applied \textbf{multi-variant clustering analysis and dimensionality reduction} to develop candidate biomarkers for disease classification and progression prediction.
  \end{itemize}
  \entrygap
  \resumeonly{\Needspace{6\baselineskip}}
  \leftright{\textbf{HPC Research Assistant (NSF-funded)}, Office of Information Technology, UT Dallas}{Aug 2023 -- May 2024}
  \begin{itemize}
    \item Deployed a brain simulation platform on a shared server with a \textbf{PostgreSQL backend} for multi-user simulation and data management.
    \item Built a \textbf{CUDA-enabled container} and integrated GPU-accelerated fiber-orientation estimation with \textbf{Slurm job scripts} into a validated UK Biobank MRI pipeline. Achieved \textbf{approximately 100\texttimes{} speedup} for fiber-orientation estimation compared with the CPU implementation.
  \end{itemize}}

%% ------------------------------------------------------------ publications
\newcommand{\PublicationsSection}{%
  %% \cvonly{\item ...} entries are left out of the resume's selected list.
  \cvsection{\resumeonly{Selected Publications}\cvonly{Publications}}
  143 citations (\link{https://scholar.google.com/citations?user=bgnVj8wAAAAJ}{Google Scholar}); research areas: \textbf{machine learning, representational drift, and whole-brain modeling}.
  \par\addvspace{4pt}\textbf{Peer-reviewed journal articles}
  \begin{itemize}
    \item \me{Wang, Y.}, Lu, Z., \& Lakshminarasimhan, K. (2026). Coordinated drift of image representations in human EEG. \status{under review}
    \item Lu, Z., \me{Wang, Y.}, \& Golomb, J. D. (2026). Achieving more human brain-like vision via human EEG representational alignment. \venue{Communications Biology}, 9, 463. \link{https://doi.org/10.1038/s42003-026-09685-w}{doi:10.1038/s42003-026-09685-w}
    \item Lu, Z., \& \me{Wang, Y.} (2025). Teaching CORnet human fMRI representations for enhanced model--brain alignment. \venue{Cognitive Neurodynamics}, 19(1), 61.
    \cvonly{%
    \item Li, Z., Huang, G., Li, Z., Li, S., \me{Wang, Y.}, Zhao, J., \ldots\ \& Zou, L. (2020). Chemosensory anhedonia in patients with schizophrenia and individuals with schizotypy: A questionnaire study. \venue{Frontiers in Psychiatry}, 11, 481.
    \item Zhao, J., Gao, Z.*, Li, Y.*, \me{Wang, Y.*}, Zhang, X., \& Zou, L. (2019). The food neophobia scale (FNS): Exploration and confirmation of factor structure in a healthy Chinese sample. \venue{Food Quality and Preference}, 79, 103791.
    \item Zhao, J.*, \me{Wang, Y.*}, Ma, Q., Zhao, J., Zhang, X., \& Zou, L. (2019). The Chemosensory Pleasure Scale: A new assessment for measuring hedonic smell and taste capacities. \venue{Chemical Senses}, 44(7), 457--464.
    \item \me{Wang, Y.}, \& Zou, L. (2018). The neural mechanisms of social pain. \venue{Progress in Biochemistry and Biophysics}, 45(7), 714--722. [in Chinese]}
  \end{itemize}
  \cvonly{\par\addvspace{2pt}*Equal contribution.}
  \par\addvspace{4pt}\textbf{Conference papers and presentations}
  \begin{itemize}
    \item \me{Wang, Y.}, \& Lakshminarasimhan, K. (2026). An efficient coding account of representational drift. \venue{International Conference on Learning Representations (ICLR 2027)}. \status{under review}
    \item Hosseinpourkhoshkbari, R., Huang, W., \me{Wang, Y.}, Helabad, A., Chen, Z., \& Golden, R. (2026). Human-likeness is multi-dimensional: Evaluating LLM-generated virtual patients against simulated patient dialogue. \venue{Conference on Language Modeling (COLM 2026)}. \status{accepted}
    \item \me{Wang, Y.}, \& Lakshminarasimhan, K. (2026). The geometry of drifting spatial representations. \venue{Cognitive Computational Neuroscience (CCN 2026)}. \status{accepted poster}
    \cvonly{%
    \item \me{Wang, Y.}, \& Lakshminarasimhan, K. (2026). When drift helps or hurts: Heterogenous effects of representational drift on decoding. \venue{Summit in BrainHealth, UT Dallas}.}
    \item Lu, Z., \me{Wang, Y.}, \& Golomb, J. D. (2024). ReAlnet: Achieving more human brain-like vision via human neural representational alignment. \venue{ICLR 2024 Workshop}. \status{accepted}
  \end{itemize}}

%% ------------------------------------------------------------ projects
\newcommand{\ProjectsSection}{%
  \cvsection{Open-Source Software}
  \leftright{\textbf{LLM-for-Zotero}: Open-source AI research agent (\link{https://github.com/yilewang/llm-for-zotero}{GitHub})}{\textit{TypeScript, 2026 -- present}}
  \subentry{Creator, lead architect, and maintainer}
  \begin{itemize}
    \item \textbf{Product ownership:} Created and shipped an AI research agent with \textbf{3,100+ GitHub stars and 650K+ plugin downloads}. Lead system architecture, implementation, \textbf{testing}, and ongoing maintenance.
    \item \textbf{Agent architecture:} Designed an \textbf{adaptive agent loop} that coordinates planning, tool execution, and continuation based on observed progress. Built an \textbf{agent harness} for instruction assembly, tool access, and runtime guardrails, including context budgets, input validation, and repeated-failure detection.
    \item \textbf{RAG and context engineering:} Combined \textbf{retrieval-augmented generation (RAG) and agentic search} for research across large document libraries. Optimized retrieval through library indexing, hybrid keyword and semantic ranking, and evidence caching. Designed automatic \textbf{context compaction} and source-coverage tracking to preserve reusable evidence across long research sessions.
    \item \textbf{Execution reliability:} Designed \textbf{state machines and action contracts} with explicit \textbf{component ownership}. The Zotero host application controls permissions, execution, and verification of saved results. Implemented \textbf{asynchronous job orchestration for batch processing}, with persisted checkpoints, cancellation, and recovery from interrupted runs. Supported undo for saved library changes.
    \item \textbf{Model and agent integration:} Integrated cloud and local models, the \textbf{Claude Agent SDK} through a companion bridge, and \textbf{Codex app-server}. Built a \textbf{Model Context Protocol (MCP) server} exposing Zotero tools to external agents. Optimized prompt caching through stable prompt prefixes and managed session reuse for continued conversations.
    \item \textbf{Testing and delivery:} Built \textbf{workflow test suites and behavioral evaluation suites} covering tool selection, permission handling, evidence quality, and verification of saved data inside Zotero. Designed automated testing, build, and publishing workflows with \textbf{GitHub Actions}, supporting \textbf{100+ releases}.
  \end{itemize}
  \cvonly{\entrygap
  \leftright{\textbf{Dynamical Network Analysis \& Statistical Toolbox} (\link{https://github.com/yilewang/dynamics_net_sys_class/tree/main}{GitHub})}{\textit{Python}}
  \begin{itemize}
    \item Teaching toolbox and notebooks for dynamical networks: SIR epidemics, network diffusion, and Erdős--Rényi graphs, linking analytic models to simulation and emergent dynamics.
  \end{itemize}}}

%% ------------------------------------------------------------ honors
\newcommand{\HonorsSection}{%
  \cvsection{Honors \& Awards}
  \leftright{Marlane Miller Visionary New Scientist Award\cvonly{ (\$25,000)}}{Oct 2023}
  \leftright{Society for Neuroscience Trainee Professional Development Award\cvonly{ (\$1,000)}}{Nov 2022}
  \leftright{UT Dallas Master's Research Fellowship\cvonly{ (\$1,500)}}{Jul 2021}}

%% ------------------------------------------------------------ skills
\newcommand{\SkillsSection}{%
  \cvsection{Skills}
  \textbf{Programming \& data:} Python, pandas, TypeScript, R, MATLAB, Bash\par
  \textbf{AI agents \& retrieval:} Adaptive agent loops, agent harnesses, retrieval-augmented generation (RAG), agentic search, context compaction, prompt caching\par
  \textbf{Software architecture \& integration:} State machines, asynchronous workflows, action contracts, recovery mechanisms, Model Context Protocol (MCP), Claude Agent SDK, Codex app-server\par
  \textbf{Machine learning \& statistics:} PyTorch, convolutional neural networks, contrastive learning, dimensionality reduction, discriminant analysis, cross-validation\par
  \textbf{Computational modeling:} Dynamical network simulation, biologically constrained neural networks, mean-field modeling, criticality analysis, information theory, Markov chain Monte Carlo (MCMC)\par
  \textbf{Testing \& evaluation:} Test-driven development (TDD), workflow test suites, regression testing\par
  \textbf{Infrastructure:} SQLite, PostgreSQL, Docker, CUDA, Slurm, high-performance computing, Git, GitHub Actions, AWS, Google Cloud\par
  \cvonly{\textbf{Neuroscience methods:} EEG/fMRI decoding, representational similarity analysis, representational drift analysis, structural MRI/DTI, individualized brain-network simulation, Brain-Score\par}}
